\documentclass[11pt]{article}

\usepackage[utf8]{inputenc}
\usepackage[T1]{fontenc}
\usepackage{amsmath,amssymb,amsthm}
\usepackage{booktabs}
\usepackage{enumitem}
\usepackage{geometry}
\usepackage{hyperref}
\usepackage[most]{tcolorbox}

\geometry{margin=1in}
\hypersetup{colorlinks=true,linkcolor=blue,urlcolor=blue}

\newcommand{\R}{\mathbb{R}}
\newcommand{\E}{\mathbb{E}}
\newcommand{\Var}{\operatorname{Var}}
\newcommand{\dd}{\,\mathrm{d}}
\newcommand{\norm}[1]{\left\lVert #1\right\rVert}

\theoremstyle{definition}
\newtheorem{definition}{Definition}[section]
\theoremstyle{plain}
\newtheorem{theorem}[definition]{Theorem}
\newtheorem{proposition}[definition]{Proposition}
\newtheorem{corollary}[definition]{Corollary}
\theoremstyle{remark}
\newtheorem{remark}[definition]{Remark}

\definecolor{PropositionBlue}{RGB}{31,78,121}
\newtcbtheorem[number within=section]{keyproposition}{Proposition}{
  enhanced,
  breakable,
  colback=PropositionBlue!4,
  colframe=PropositionBlue,
  colbacktitle=PropositionBlue,
  coltitle=white,
  fonttitle=\bfseries,
  boxrule=0.9pt,
  arc=1.5mm,
  left=2mm,right=2mm,top=1.5mm,bottom=1.5mm
}{prop}

\tcolorboxenvironment{theorem}{
  enhanced,breakable,
  colback=PropositionBlue!4,colframe=PropositionBlue,
  boxrule=0.9pt,arc=1.5mm,
  left=2mm,right=2mm,top=1.5mm,bottom=1.5mm
}
\tcolorboxenvironment{proposition}{
  enhanced,breakable,
  colback=PropositionBlue!4,colframe=PropositionBlue,
  boxrule=0.9pt,arc=1.5mm,
  left=2mm,right=2mm,top=1.5mm,bottom=1.5mm
}

\title{Summer Bootcamp -- Session 2 Notes\\Calculus and Ordinary Differential Equations}
\author{Nathan Sauldubois}
\date{August 2026}

\begin{document}
\maketitle

\section*{Purpose of the Session}

This second three-hour session reviews the calculus tools needed for financial
models, numerical methods, and optimization. The first part covers differential
and integral calculus. The second part develops scalar and multidimensional
linear ordinary differential equations, their numerical resolution, and their
applications to risk-free growth, continuous discounting, and the Vasicek
short-rate model.

\section{Landau Notation and Local Approximation}

Landau notation compares the size of two quantities near a limit point. Let
$f$ and $g$ be defined near $a$, with $g(x)\neq0$ for $x\neq a$.

\begin{definition}[Big $O$ and small $o$ in one dimension]
We write $f(x)=O(g(x))$ as $x\to a$ if there exist $C,\delta>0$ such that
\[
0<|x-a|<\delta\quad\Longrightarrow\quad |f(x)|\leq C|g(x)|.
\]
We write $f(x)=o(g(x))$ as $x\to a$ if
\[
\frac{f(x)}{g(x)}\longrightarrow0.
\]
\end{definition}

Thus $O(g)$ means ``at most the scale of $g$'', whereas $o(g)$ means
``negligible compared with $g$''. For example, as $h\to0$,
\[
h^3=o(h^2),\qquad 3h^2-h^3=O(h^2),\qquad
\sin h=h+O(h^3).
\]
Note that $h^2=O(h^2)$ but $h^2\neq o(h^2)$.

For maps $F:\R^d\to\R^m$ and $G:\R^d\to[0,\infty)$, the same definitions
use norms:
\[
F(h)=O(G(h))\iff \norm{F(h)}\leq CG(h)\quad\text{near }0,
\]
\[
F(h)=o(G(h))\iff \frac{\norm{F(h)}}{G(h)}\longrightarrow0.
\]
For $h=(h_1,h_2)$, one has
\[
h_1h_2=O(\norm h^2),
\qquad h_1^2h_2=o(\norm h^2).
\]
These notations encode the remainders in differentiability and Taylor
expansions throughout the session.

\section{One-Dimensional Differentiation}

For $f:\R\to\R$, the derivative at $x$ is
\[
f'(x)=\lim_{h\to0}\frac{f(x+h)-f(x)}{h},
\]
when this limit exists. It measures the instantaneous rate of change of $f$ and
the slope of the tangent line at $(x,f(x))$.

\subsection{Derivative and Local Linear Approximation}

If $h$ is small, then
\[
f(x+h)\approx f(x)+f'(x)h.
\]
This is the first example of a central idea in applied mathematics: replace a
complicated nonlinear object by a simpler local approximation.

More precisely, differentiability at $x$ means that there is a remainder
$r(h)$ such that
\[
f(x+h)=f(x)+f'(x)h+r(h),
\qquad
\frac{r(h)}{|h|}\longrightarrow0.
\]
This formulation, rather than the difference quotient alone, generalizes
directly to several variables.

\begin{definition}[$C^1$ regularity in one dimension]
Let $I\subset\R$ be an open interval. A function $f:I\to\R$ belongs to
$C^1(I)$ if it is differentiable on $I$ and its derivative $f'$ is continuous
on $I$.
\end{definition}

Differentiability alone does not imply $C^1$ regularity. For example,
$x\mapsto e^x$ and $x\mapsto|x|^3$ are in $C^1(\R)$, whereas
$x\mapsto|x|$ is not differentiable at the origin.

\subsection{One-Dimensional Derivative Toolkit}

For differentiable functions $f,g$ and constants $a,b$,
\[
(af+bg)'=af'+bg',
\qquad
(fg)'=f'g+fg',
\]
\[
\left(\frac{f}{g}\right)'=\frac{f'g-fg'}{g^2},
\qquad g(x)\neq0.
\]
\begin{keyproposition}{Chain rule in one dimension}{chain-rule-1d}
Let $g:I\to J$ be differentiable at $x$ and let $f:J\to\R$ be
differentiable at $g(x)$. Then $f\circ g$ is differentiable at $x$ and
\[
\boxed{(f\circ g)'(x)=f'(g(x))\,g'(x).}
\]
The outer derivative is evaluated at the inner function, then multiplied by
the derivative of the inner function.
\end{keyproposition}

For example,
\[
\frac{\dd}{\dd x}(2x+1)^4=4(2x+1)^3\cdot2=8(2x+1)^3.
\]

The complete one-dimensional toolkit used throughout the course is therefore
\[
\boxed{
\begin{aligned}
(af+bg)'&=af'+bg',
& (fg)'&=f'g+fg',\\
\left(\frac fg\right)'&=\frac{f'g-fg'}{g^2},
& (f\circ g)'&=(f'\circ g)g',\\
(g^\alpha)'&=\alpha g^{\alpha-1}g',
& (e^g)'&=e^g g',\\
(\log g)'&=\frac{g'}{g}.
\end{aligned}}
\]
The elementary derivatives most frequently used in this course are summarized
below:
\[
\begin{array}{c|c@{\qquad}c|c}
f(x)&f'(x)&f(x)&f'(x)\\ \hline
x^\alpha&\alpha x^{\alpha-1}&e^{ax}&ae^{ax}\\
\log x&1/x&\log(g(x))&g'(x)/g(x)\\
\sin x&\cos x&\cos x&-\sin x\\
a^x&a^x\log a&g(x)^\alpha&\alpha g(x)^{\alpha-1}g'(x)
\end{array}
\]
The power formulas are understood on domains where the expressions are
well-defined.

\subsection{Rolle's Theorem and the Mean Value Theorem}

\begin{theorem}[Rolle's theorem]
If $f$ is continuous on $[a,b]$, differentiable on $(a,b)$, and
$f(a)=f(b)$, then there exists $c\in(a,b)$ such that
\[
f'(c)=0.
\]
\end{theorem}

\begin{proof}
By continuity, $f$ attains a minimum and a maximum on $[a,b]$. If both are
equal, $f$ is constant and every interior point works. Otherwise, at least one
extremum is attained at an interior point because $f(a)=f(b)$. Fermat's
theorem then gives $f'(c)=0$ at that interior extremum.
\end{proof}

Geometrically, a differentiable curve that starts and ends at the same height
must have a horizontal tangent somewhere in between. For example,
$f(x)=x^2-4x$ satisfies $f(0)=f(4)=0$, and Rolle's point is $c=2$.

\begin{theorem}[Mean value theorem]
If $f$ is continuous on $[a,b]$ and differentiable on $(a,b)$, then there
exists $c\in(a,b)$ such that
\[
f(b)-f(a)=f'(c)(b-a).
\]
\end{theorem}

\begin{proof}
Let $\ell(x)$ be the affine function joining $(a,f(a))$ and $(b,f(b))$, and set
$g=f-\ell$. Then $g(a)=g(b)=0$. Rolle's theorem gives a point $c$ with
$g'(c)=0$. Since
\[
\ell'(x)=\frac{f(b)-f(a)}{b-a},
\]
the equality $g'(c)=0$ is exactly the desired formula.
\end{proof}

\begin{keyproposition}{Monotonicity consequences of the mean value theorem}{mvt-monotonicity}
Let $f$ be differentiable on an interval $I$. Then
\[
\boxed{
\begin{aligned}
f'\geq0\text{ on }I&\Longrightarrow f\text{ is nondecreasing},\\
f'>0\text{ on }I&\Longrightarrow f\text{ is strictly increasing},\\
f'\leq0\text{ on }I&\Longrightarrow f\text{ is nonincreasing},\\
f'=0\text{ on }I&\Longrightarrow f\text{ is constant}.
\end{aligned}}
\]
The analogous strict condition $f'<0$ implies that $f$ is strictly decreasing.
\end{keyproposition}

Indeed, if $x<y$, the mean value theorem gives $c\in(x,y)$ such that
\[
f(y)-f(x)=f'(c)(y-x).
\]
Since $y-x>0$, the sign of the increment is controlled by the sign of the
derivative.

\begin{keyproposition}{Derivative bound and uniqueness}{mvt-lipschitz}
If $|f'(x)|\leq L$ for every $x$ in an interval $I$, then
\[
\boxed{|f(y)-f(x)|\leq L|y-x|\qquad\text{for all }x,y\in I.}
\]
Thus $f$ is $L$-Lipschitz on $I$. Moreover, if $f'>0$ everywhere or $f'<0$
everywhere, then $f$ is one-to-one; in particular,
\[
\boxed{f(x)=0\text{ has at most one solution in }I.}
\]
\end{keyproposition}

\section{Multivariate Calculus}

\subsection{Differentiability, Partial Derivatives, and Gradient}

For $f:\R^n\to\R$, the partial derivative with respect to $x_i$ is
\[
\partial_{x_i}f(x)
=
\lim_{h\to0}
\frac{f(x_1,\ldots,x_i+h,\ldots,x_n)-f(x)}{h}.
\]
The gradient is the vector of partial derivatives:
\[
\nabla f(x)=
\begin{pmatrix}
\partial_{x_1}f(x)\\
\vdots\\
\partial_{x_n}f(x)
\end{pmatrix}.
\]
The directional derivative in a unit direction $u$ is
\[
D_uf(x)=\nabla f(x)^\top u.
\]
The gradient points in the direction of steepest ascent, and its norm is the
maximal instantaneous rate of increase.

The existence of all partial derivatives is not, by itself, enough to ensure a
good linear approximation. The correct multivariate definition is the
following.

\begin{definition}[Fréchet differentiability]
A map $g:\R^n\to\R^m$ is differentiable at $x$ if there is a linear map
$Dg(x):\R^n\to\R^m$ such that
\[
g(x+h)=g(x)+Dg(x)h+r(h),
\qquad
\frac{\norm{r(h)}}{\norm{h}}\longrightarrow0.
\]
In canonical coordinates, the matrix of $Dg(x)$ is the Jacobian $J_g(x)$.
For scalar-valued $f$, $Df(x)h=\nabla f(x)^\top h$.
\end{definition}

\begin{definition}[$C^1$ regularity in several dimensions]
Let $U\subset\R^n$ be open. A map $g:U\to\R^m$ belongs to
$C^1(U;\R^m)$ if it is differentiable on $U$ and
\[
x\longmapsto Dg(x)\in\mathcal L(\R^n,\R^m)
\]
is continuous. Equivalently, all first-order partial derivatives
$\partial_{x_j}g_i$ exist and are continuous on $U$.
\end{definition}

The word \emph{linear} is essential. Once the base point $x$ is fixed, the
differential acts linearly on increments: for all $h,k\in\R^n$ and scalars
$\alpha,\beta$,
\[
\boxed{
Dg(x)(\alpha h+\beta k)
=\alpha Dg(x)h+\beta Dg(x)k.}
\]
Equivalently, because $Dg(x)$ is represented by the Jacobian,
\[
J_g(x)(\alpha h+\beta k)
=\alpha J_g(x)h+\beta J_g(x)k.
\]
The original map $g$ need not be linear. Differentiability means that its
first-order variation around the fixed point $x$ is described by this linear
map. The derivative is also unique: two linear maps satisfying the defining
remainder condition must coincide.

If the partial derivatives exist and are continuous in a neighborhood of $x$,
then the function is differentiable at $x$. This is the regularity criterion
used throughout the course.

\subsection{Computing a Gradient: Partials and Increments}

The differential of $g$ at $x$ is the linear map
\[
\dd g_x(h)=Dg(x)h=J_g(x)h.
\]
For a scalar function, this becomes
\[
\dd f_x(h)=\nabla f(x)^\top h.
\]
The notation emphasizes that the differential is linear in the increment
$h$, even when the original function is nonlinear.

Consider the example used in the slides:
\[
f(x,y)=x^2+3xy+2y^2.
\]
The coordinate calculation gives
\[
\partial_xf=2x+3y,
\qquad
\partial_yf=3x+4y,
\]
and hence
\[
\boxed{\nabla f(x,y)=\begin{pmatrix}2x+3y\\3x+4y\end{pmatrix}.}
\]

\subsection{Core Gradient and Jacobian Formulas}

For constant vectors $a,b$, matrices $A,Q$, and a differentiable scalar
function $\phi$, the following formulas are fundamental:
\[
\begin{aligned}
\nabla(a^\top x)&=a,\\
\nabla(b^\top Ax)&=A^\top b,\\
\nabla\!\left(\tfrac12\norm{Ax-b}^2\right)&=A^\top(Ax-b),\\
\nabla(x^\top Qx)&=(Q+Q^\top)x,\\
\nabla\phi(a^\top x)&=\phi'(a^\top x)a.
\end{aligned}
\]

\subsubsection*{Scalar product: two calculations}

Let $p(x,y)=x^\top y$ for $x,y\in\R^n$. Coordinatewise,
\[
\partial_{x_i}p=y_i,
\qquad
\partial_{y_i}p=x_i,
\qquad
\nabla p(x,y)=\begin{pmatrix}y\\x\end{pmatrix}.
\]
For increments $(h,k)$,
\[
p(x+h,y+k)-p(x,y)
=y^\top h+x^\top k+h^\top k.
\]
The last term is quadratic in the increments. Therefore
\[
\boxed{Dp(x,y)(h,k)=y^\top h+x^\top k.}
\]
The coordinate calculation and the increment calculation identify the same
linear map.

\subsubsection*{Quadratic form: two calculations}

Let
\[
f(x)=x^\top Qx,
\qquad Q\in\R^{n\times n}.
\]
Writing the expression in coordinates and differentiating gives
\[
\nabla f(x)=(Q+Q^\top)x.
\]
The increment calculation gives
\[
f(x+h)-f(x)=x^\top Qh+h^\top Qx+h^\top Qh.
\]
Hence
\[
Df(x)h=((Q+Q^\top)x)^\top h,
\]
which recovers the same gradient. If $Q$ is symmetric, then
$\nabla f(x)=2Qx$.
This formula is important in
least squares, portfolio variance, and quadratic optimization.

\subsubsection*{Jacobian of a vector-valued map}

For a vector-valued function $g:\R^n\to\R^m$,
\[
g(x)=
\begin{pmatrix}
g_1(x)\\
\vdots\\
g_m(x)
\end{pmatrix},
\]
the Jacobian matrix is
\[
J_g(x)=
\begin{pmatrix}
\partial_{x_1}g_1(x)&\cdots&\partial_{x_n}g_1(x)\\
\vdots&\ddots&\vdots\\
\partial_{x_1}g_m(x)&\cdots&\partial_{x_n}g_m(x)
\end{pmatrix}.
\]
It is the matrix of the best local linear approximation of $g$.

\subsubsection*{Jacobian example}

For
\[
\Phi(u,v)=(ue^v,u^2+v),
\]
we have
\[
J_\Phi(u,v)=
\begin{pmatrix}
e^v&ue^v\\
2u&1
\end{pmatrix},
\qquad
\det J_\Phi(u,v)=e^v(1-2u^2).
\]
The determinant measures local area distortion. When it is close to zero, the
mapping is locally close to singular.

\clearpage
\subsection{Multivariate Chain Rule}

\begin{keyproposition}{Multivariate chain rule}{chain-rule-multivariate}
Let $g:\R^n\to\R^m$ be differentiable at $x$ and let
$f:\R^m\to\R^p$ be differentiable at $g(x)$. Then $f\circ g$ is
differentiable at $x$ and
\[
\boxed{J_{f\circ g}(x)=J_f(g(x))\,J_g(x).}
\]
For a scalar-valued outer function $f:\R^m\to\R$, this becomes
\[
\boxed{\nabla(f\circ g)(x)=J_g(x)^\top\nabla f(g(x)).}
\]
\end{keyproposition}

This form is used constantly in sensitivity analysis, automatic
differentiation, and optimization.

\begin{proof}[Proof of the chain rule]
Write
\[
g(x+h)=g(x)+J_g(x)h+r_g(h),
\qquad \norm{r_g(h)}=o(\norm h).
\]
Apply the differentiability expansion of $f$ at $g(x)$ to the increment
$J_g(x)h+r_g(h)$. The linear term is
$J_f(g(x))J_g(x)h$; all remaining terms are $o(\norm h)$. The coefficient of
the first-order term is therefore the product of the two Jacobians.
\end{proof}

The same formula can be read coordinatewise. Writing
$g=(g_1,\ldots,g_m)$, for every $j$,
\[
\boxed{
\partial_{x_j}(f\circ g)(x)
=\sum_{i=1}^m
\partial_{y_i}f(g(x))\,\partial_{x_j}g_i(x).}
\]
This is the $j$-th component of
$J_g(x)^\top\nabla f(g(x))$.

\section{Implicit and Inverse Function Theorems}

\subsection{Implicit Function Theorem}

Suppose a relation $F(x,y)=0$ defines $y$ indirectly as a function of $x$.
The implicit function theorem gives a local existence criterion and a formula
for the derivative without requiring an explicit solution.

\begin{theorem}[Implicit function theorem, scalar form]
Let $F:\R^2\to\R$ be continuously differentiable near $(x_0,y_0)$, with
\[
F(x_0,y_0)=0,
\qquad
\partial_yF(x_0,y_0)\neq0.
\]
Then, near $x_0$, there is a unique differentiable function $y=\varphi(x)$
such that $\varphi(x_0)=y_0$ and
\[
F(x,\varphi(x))=0.
\]
Moreover,
\[
\boxed{\varphi'(x)=-\frac{\partial_xF(x,\varphi(x))}
{\partial_yF(x,\varphi(x))}}.
\]
\end{theorem}

The existence statement is deeper than the derivative formula. Once existence
is known, the formula follows immediately by differentiating
$F(x,\varphi(x))=0$ and applying the chain rule:
\[
0=\partial_xF+\partial_yF\,\varphi'.
\]

\subsection{Implicit Differentiation on the Unit Circle}

For the unit circle $F(x,y)=x^2+y^2-1$, every point with $y\neq0$ satisfies
the condition $\partial_yF=2y\neq0$. Hence the circle is locally a graph and
\[
\frac{\dd y}{\dd x}=-\frac{x}{y}.
\]
At $(1,0)$ the condition fails: the tangent is vertical, so $y$ cannot be a
single differentiable function of $x$ near that point.

\subsection{Inverse Function Theorem and Newton's Method}

\begin{theorem}[Inverse function theorem]
Let $g:\R^n\to\R^n$ be continuously differentiable near $x_0$. If
$J_g(x_0)$ is invertible, then $g$ has a differentiable local inverse near
$g(x_0)$, and
\[
J_{g^{-1}}(g(x_0))=J_g(x_0)^{-1}.
\]
\end{theorem}

The inverse problem starts from a target $y\in\R^n$ and asks for $x$ such that
\[
\boxed{g(x)=y.}
\]
Set $F(x)=g(x)-y$. At an iterate $x_k$, linearization gives
\[
F(x_k+s)\approx F(x_k)+J_g(x_k)s.
\]
Imposing the linearized equation yields Newton's algorithm:
\[
\boxed{J_g(x_k)s_k=y-g(x_k)},
\qquad
\boxed{x_{k+1}=x_k+s_k}.
\]
One solves this linear system directly and never forms $J_g(x_k)^{-1}$.
Newton converges rapidly near a regular solution, but a poor initial guess or
a nearly singular Jacobian can cause failure.

In one dimension, bisection provides a robust alternative. If $F$ is
continuous and $F(a)F(b)<0$, repeat the following steps:
\begin{enumerate}
\item compute $m=(a+b)/2$;
\item if $F(a)F(m)\leq0$, replace $b$ by $m$;
\item otherwise, replace $a$ by $m$;
\item stop when the interval length or the residual is sufficiently small.
\end{enumerate}
After $k$ iterations,
\[
|x_k-x^\star|\leq\frac{b-a}{2^{k+1}}.
\]
Bisection is guaranteed but only linearly convergent; Newton is locally
quadratic but needs a derivative and a good initial point.

\subsubsection*{Implied volatility as a function of strike and maturity}

For a call with spot $S$, strike $K$, maturity $T$, constant rate $r$, and
volatility $\sigma$, the Black--Scholes price is
\[
C_{\mathrm{BS}}(S,K,T,r,\sigma)
=S\Phi(d_1)-Ke^{-rT}\Phi(d_2),
\]
where
\[
d_1=\frac{\log(S/K)+(r+\tfrac12\sigma^2)T}{\sigma\sqrt T},
\qquad
d_2=d_1-\sigma\sqrt T.
\]
For every quoted pair $(K,T)$, implied volatility is defined by
\[
\boxed{C_{\mathrm{BS}}(S,K,T,r,\sigma_{\mathrm{imp}}(K,T))
=C^{\mathrm{mkt}}(K,T).}
\]
Define
\[
F(K,T,\sigma)
=C_{\mathrm{BS}}(S,K,T,r,\sigma)-C^{\mathrm{mkt}}(K,T).
\]
Since
\[
\partial_\sigma F=\mathrm{Vega}>0
\]
for $T>0$ in the nondegenerate Black--Scholes setting, the implicit function
theorem applies. It provides a locally unique differentiable function
$(K,T)\mapsto\sigma_{\mathrm{imp}}(K,T)$ satisfying $F=0$.
The map $K\mapsto\sigma_{\mathrm{imp}}(K,T)$ is generally not flat in market
data. A U-shaped cross-section is called a volatility smile; a systematically
downward-sloping cross-section is called a volatility skew.

\section{Second Derivatives and Hessians}

\subsection{Hessian and Symmetry of Mixed Derivatives}

For $f:\R^n\to\R$, the Hessian matrix is
\[
H_f(x)=
\begin{pmatrix}
\partial_{x_1x_1}^2f(x)&\cdots&\partial_{x_1x_n}^2f(x)\\
\vdots&\ddots&\vdots\\
\partial_{x_nx_1}^2f(x)&\cdots&\partial_{x_nx_n}^2f(x)
\end{pmatrix}.
\]
If $f$ is sufficiently regular, Schwarz theorem gives
\[
\partial_{x_ix_j}^2f=\partial_{x_jx_i}^2f,
\]
so the Hessian is symmetric.

\begin{theorem}[Schwarz--Clairaut theorem]
If the second partial derivatives of $f$ are continuous in a neighborhood of
$x$, then
\[
\partial_{x_ix_j}^2f(x)=\partial_{x_jx_i}^2f(x).
\]
Consequently $H_f(x)$ is symmetric.
\end{theorem}

\begin{proof}[Proof idea in two dimensions]
Apply the fundamental theorem of calculus twice to the rectangular increment
\[
f(x+h,y+k)-f(x+h,y)-f(x,y+k)+f(x,y).
\]
Integrating first in $x$ and then in $y$ represents it as the double integral
of $f_{xy}$; reversing the order represents the same increment using $f_{yx}$.
Continuity and shrinking the rectangle to $(x,y)$ give $f_{xy}=f_{yx}$.
\end{proof}

\subsection{Convexity and Second-Order Classification}

\begin{proposition}[Hessian criterion for convexity]
Let $f$ be twice continuously differentiable on a convex open set. If
$H_f(x)$ is positive semidefinite for every $x$, then $f$ is convex. If it is
positive definite everywhere, then $f$ is strictly convex.
\end{proposition}

\begin{proof}
Fix $x,y$ and define the one-dimensional function
$g(t)=f(x+t(y-x))$ for $t\in[0,1]$. The chain rule gives
\[
g''(t)=(y-x)^\top H_f(x+t(y-x))(y-x)\geq0.
\]
Thus $g$ is convex on every line segment, which is exactly convexity of $f$.
\end{proof}

At a critical point $x^\star$, a positive definite Hessian implies a strict
local minimum, a negative definite Hessian implies a strict local maximum, and
an indefinite Hessian implies a saddle point.

\subsubsection*{Worked example}

Let
\[
f(x,y)=x^4+2x^2y^2+y^4=(x^2+y^2)^2.
\]
Then
\[
\nabla f(x,y)=
\begin{pmatrix}
4x(x^2+y^2)\\
4y(x^2+y^2)
\end{pmatrix}
\]
and
\[
H_f(x,y)=
\begin{pmatrix}
12x^2+4y^2&8xy\\
8xy&4x^2+12y^2
\end{pmatrix}.
\]
This Hessian is positive semidefinite, which is consistent with the convexity of
$f$.

\subsection{Hands-On Example: Least Squares}

For data matrix $A\in\R^{m\times n}$ and observations $b\in\R^m$, define
\[
L(x)=\frac12\norm{Ax-b}^2.
\]
Then
\[
\nabla L(x)=A^\top(Ax-b),
\qquad
H_L(x)=A^\top A.
\]
Critical points solve the normal equations
\[
A^\top Ax=A^\top b.
\]
For
\[
A=\begin{pmatrix}1&0\\1&1\\1&2\end{pmatrix},
\qquad
b=\begin{pmatrix}1\\2\\2\end{pmatrix},
\]
one obtains
\[
A^\top A=\begin{pmatrix}3&3\\3&5\end{pmatrix},
\qquad
A^\top b=\begin{pmatrix}5\\6\end{pmatrix},
\qquad
x^\star=\begin{pmatrix}7/6\\1/2\end{pmatrix}.
\]
Because $A$ has full column rank, $A^\top A$ is positive definite, so this
critical point is the unique global minimizer.

\section{Taylor Expansion}

\subsection{One-Dimensional Taylor Theorem}

If $f$ is sufficiently smooth near $x$, then
\[
f(x+h)=f(x)+f'(x)h+\frac{1}{2}f''(x)h^2+\cdots+
\frac{1}{k!}f^{(k)}(x)h^k+R_k(h).
\]
The most important approximations for this bootcamp are
\[
f(x+h)\approx f(x)+f'(x)h
\]
and
\[
f(x+h)\approx f(x)+f'(x)h+\frac{1}{2}f''(x)h^2.
\]

\begin{theorem}[Taylor theorem with Lagrange remainder]
If $f$ is $(k+1)$ times continuously differentiable between $x$ and $x+h$,
then for some $\xi$ between these two points,
\[
f(x+h)=\sum_{j=0}^k\frac{f^{(j)}(x)}{j!}h^j
+\frac{f^{(k+1)}(\xi)}{(k+1)!}h^{k+1}.
\]
\end{theorem}

The theorem gives an error bound. If $|f^{(k+1)}|\leq M$ on the interval, then
\[
|R_k(h)|\leq\frac{M}{(k+1)!}|h|^{k+1}.
\]

\subsection{Standard Expansions and Error Control}

\[
e^x=1+x+\frac{x^2}{2}+\frac{x^3}{6}+\cdots,
\]
\[
\log(1+x)=x-\frac{x^2}{2}+\frac{x^3}{3}-\cdots,
\]
\[
\sin x=x-\frac{x^3}{6}+\frac{x^5}{120}-\cdots.
\]

For instance, the second-order approximation $e^{0.1}\approx1+0.1+0.1^2/2$
gives $1.105$. Since the third derivative of $e^x$ is at most $e^{0.1}$ on
$[0,0.1]$, the error is bounded by
\[
\frac{e^{0.1}}{6}(0.1)^3<1.85\times10^{-4}.
\]

\subsection{Multivariate Taylor Expansion and Critical Points}

\begin{keyproposition}{Second-order multivariate Taylor expansion}{taylor-order-two}
Let $f:\R^n\to\R$ be twice continuously differentiable in a neighborhood
of $x$. As $h\to0$,
\[
\boxed{
f(x+h)=f(x)+\nabla f(x)^\top h
+\frac{1}{2}h^\top H_f(x)h+o(\norm h^2).}
\]
Equivalently, the remainder divided by $\norm h^2$ converges to zero. If the
third derivatives are bounded locally, the remainder is $O(\norm h^3)$.
\end{keyproposition}

The first-order term measures local sensitivity. The second-order term measures
curvature. This formula is the bridge to second-order optimization methods.

If $\nabla f(x^\star)=0$, then
\[
f(x^\star+h)\approx f(x^\star)+\frac{1}{2}h^\top H_f(x^\star)h.
\]
If $H_f(x^\star)$ is positive definite, $x^\star$ is a local minimum. If it is
negative definite, $x^\star$ is a local maximum. If it is indefinite,
$x^\star$ is a saddle point.

\section{Integration}

\subsection{Fundamental Theorem and Main Techniques}

\begin{theorem}[Fundamental theorem of calculus]
If $f$ is continuous on $[a,b]$ and
\[
F(x)=\int_a^x f(t)\dd t,
\]
then $F$ is differentiable and
\[
\boxed{F'(x)=\frac{\dd}{\dd x}\int_a^x f(t)\dd t=f(x).}
\]
Consequently, if $G'=f$, then
\[
\boxed{\int_a^bf(x)\dd x=G(b)-G(a).}
\]
\end{theorem}

\begin{proof}[Proof of the derivative statement]
For $h\neq0$,
\[
\frac{F(x+h)-F(x)}{h}=\frac1h\int_x^{x+h}f(t)\dd t.
\]
The right-hand side is the average value of $f$ over a shrinking interval.
Continuity implies that this average converges to $f(x)$.
\end{proof}

If $F'(x)=f(x)$, then
\[
\int_a^b f(x)\dd x=F(b)-F(a).
\]
For example,
\[
\int_0^T e^{rt}\dd t=
\begin{cases}
\dfrac{e^{rT}-1}{r},& r\neq0,\\
T,& r=0.
\end{cases}
\]

\subsubsection*{Change of variables}

If $u=\phi(x)$ and $\phi$ is differentiable, then
\[
\int_a^b f(\phi(x))\phi'(x)\dd x
=
\int_{\phi(a)}^{\phi(b)}f(u)\dd u.
\]
For example,
\[
\int_0^1\frac{1}{\sqrt{1-x^2}}\dd x.
\]
Use $x=\sin u$. Then $\dd x=\cos u\dd u$, and $u$ goes from $0$ to $\pi/2$.
Therefore
\[
\int_0^1\frac{1}{\sqrt{1-x^2}}\dd x
=
\int_0^{\pi/2}1\dd u
=\frac{\pi}{2}.
\]

\subsubsection*{Change of variables in arbitrary dimension}

The one-dimensional derivative is replaced by the determinant of the
Jacobian. It measures the local change of volume produced by the map.

A bijection has a unique inverse. A $C^1$ diffeomorphism is a bijection
$\Phi$ such that both $\Phi$ and $\Phi^{-1}$ are continuously differentiable.

\begin{theorem}[Change of variables in $\R^n$]
Let $U,V\subset\R^n$ be open sets and let
$\Phi:U\to V$ be a continuously differentiable bijection whose inverse is
continuously differentiable. For every integrable function $f:V\to\R$,
\[
\boxed{
\int_V f(y)\dd y
=\int_U f(\Phi(x))\left|\det J_\Phi(x)\right|\dd x.}
\]
Equivalently, for a measurable domain $D\subset U$,
\[
\boxed{
\int_{\Phi(D)} f(y)\dd y
=\int_D f(\Phi(x))\left|\det J_\Phi(x)\right|\dd x.}
\]
\end{theorem}

The absolute value is essential: integration measures unsigned volume, while
$\det J_\Phi$ is negative when the transformation reverses orientation. The
formula is the multidimensional version of
$\dd u=|\phi'(x)|\dd x$.

\begin{proof}[Geometric derivation]
Near a point $x$, differentiability gives
\[
\Phi(x+h)=\Phi(x)+J_\Phi(x)h+o(\|h\|).
\]
Thus a sufficiently small box near $x$ is transformed, to first order, by the
linear map $J_\Phi(x)$. A linear map multiplies volumes by the absolute value
of its determinant. Summing these local approximations and passing to the
limit gives the integral formula.
\end{proof}

\paragraph{Polar coordinates.}
For
\[
\Phi(r,\theta)=(r\cos\theta,r\sin\theta),
\]
the Jacobian is
\[
J_\Phi(r,\theta)=
\begin{pmatrix}
\cos\theta&-r\sin\theta\\
\sin\theta&r\cos\theta
\end{pmatrix},
\qquad
|\det J_\Phi|=r.
\]
Hence, for example,
\[
\iint_{x^2+y^2\le R^2}(x^2+y^2)\dd x\dd y
=\int_0^{2\pi}\int_0^R r^2\,r\dd r\dd\theta
=\frac{\pi R^4}{2}.
\]

\subsubsection*{Integration by parts}

\begin{keyproposition}{Integration by parts}{integration-by-parts}
Let $u,v\in C^1([a,b])$. Integrating the product rule
$(uv)'=u'v+uv'$ gives
\[
\boxed{
\int_a^b u(x)v'(x)\dd x
=\left[u(x)v(x)\right]_a^b
-\int_a^b u'(x)v(x)\dd x.}
\]
Equivalently, $\int_a^b u\dd v=[uv]_a^b-\int_a^b v\dd u$.
\end{keyproposition}

For $a>0$,
\[
\int_0^T te^{-at}\dd t
=
\frac{1-e^{-aT}(1+aT)}{a^2}.
\]

\clearpage
\subsection{Parameter-Dependent Integrals and Leibniz Rule}

\begin{keyproposition}{Differentiation under the integral sign}{differentiate-integral}
Let $F(\theta)=\int_a^b f(x,\theta)\dd x$. Assume that
$\partial_\theta f(x,\theta)$ exists near $\theta_0$ and that there is an
integrable function $g$ such that
\[
|\partial_\theta f(x,\theta)|\leq g(x)
\]
for all nearby $\theta$. Then $F$ is differentiable at $\theta_0$ and
\[
\boxed{
F'(\theta_0)=\int_a^b\partial_\theta f(x,\theta_0)\dd x.}
\]
The domination assumption justifies exchanging differentiation and integration.
\end{keyproposition}

\begin{keyproposition}{Leibniz rule with moving boundaries}{leibniz-rule}
Let
\[
F(\theta)=\int_{a(\theta)}^{b(\theta)} f(x,\theta)\dd x.
\]
If $a,b$ are differentiable, $f$ is continuous, and differentiation under the
integral sign is justified, then
\[
\boxed{
F'(\theta)
=
f(b(\theta),\theta)b'(\theta)
-f(a(\theta),\theta)a'(\theta)
+\int_{a(\theta)}^{b(\theta)}
\partial_\theta f(x,\theta)\dd x.}
\]
\end{keyproposition}

This is the main tool behind differentiating option-pricing integrals with
respect to model parameters.

\subsubsection*{Sensitivity of an expectation}

If a price or quantity can be written as
\[
P(\theta)=\E[g(X_\theta)]
=\int_\R g(x)p(x,\theta)\dd x,
\]
then, when regularity conditions justify the interchange,
\[
\frac{\partial P}{\partial\theta}
=
\int_\R g(x)\partial_\theta p(x,\theta)\dd x.
\]
This is a calculus view of Greeks.

\section{Ordinary Differential Equations}

\subsection{Theoretical Framework: Initial Value Problems and Existence}

An ordinary differential equation describes the time evolution of a function:
\[
y'(t)=f(t,y(t)),
\qquad y(0)=y_0.
\]
The derivative tells us the instantaneous direction of motion. The initial
condition selects one trajectory among all possible solutions.

\begin{keyproposition}{Picard--Lindelöf existence and uniqueness theorem}{picard-lindelof}
If $f(t,y)$ is continuous near $(t_0,y_0)$ and locally Lipschitz in $y$, then
the initial value problem
\[
\boxed{y'(t)=f(t,y(t)),\qquad y(t_0)=y_0}
\]
has a unique solution on some interval around $t_0$.
\end{keyproposition}

Local Lipschitz continuity in $y$ means that near every $(t_0,y_0)$ there
exist a neighborhood and $L<\infty$ such that
\[
|f(t,y_1)-f(t,y_2)|\leq L|y_1-y_2|
\]
for points in that neighborhood. A continuous derivative $\partial_yf$ is a
standard sufficient condition.

The continuity assumption gives existence; the Lipschitz condition prevents
two different solution curves from leaving the same initial point. For
example, $y'=\sqrt{|y|}$ with $y(0)=0$ illustrates why uniqueness can fail
when local Lipschitz continuity is absent.

\subsection{General Calculations for Scalar Linear Equations}

An equation is linear in the unknown function if it can be written
\[
y'(t)=a(t)y(t)+b(t).
\]
The term $b(t)$ is the forcing term. We first solve the homogeneous equation
\[
y'(t)=a(t)y(t),
\qquad y(0)=y_0.
\]

\subsubsection*{Fundamental solution}

Define
\[
\boxed{\Phi(t,0)=\exp\left(\int_0^t a(u)\dd u\right).}
\]
Then
\[
\partial_t\Phi(t,0)=a(t)\Phi(t,0),
\qquad \Phi(0,0)=1,
\]
and the homogeneous solution is
\[
\boxed{y(t)=\Phi(t,0)y_0.}
\]
For a constant coefficient $a(t)=k$, this reduces to
$\Phi(t,0)=e^{kt}$. Positive $k$ gives growth and negative $k$ gives decay.

\subsubsection*{Variation of the constant}

For the forced equation, seek
\[
y(t)=c(t)\Phi(t,0).
\]
Differentiating gives
\[
y'=c'\Phi+c\,a(t)\Phi.
\]
After substitution, the homogeneous terms cancel and
\[
c'(t)\Phi(t,0)=b(t).
\]
Since $c(0)=y_0$,
\[
c(t)=y_0+\int_0^t\Phi(s,0)^{-1}b(s)\dd s.
\]
Introducing
\[
\Phi(t,s)=\Phi(t,0)\Phi(s,0)^{-1}
=\exp\left(\int_s^t a(u)\dd u\right),
\]
we obtain the main solution formula.

\begin{keyproposition}{Variation of constants for a scalar linear ODE}{scalar-variation-constants}
If $a,b$ are continuous, the initial value problem
\[
y'(t)=a(t)y(t)+b(t),\qquad y(0)=y_0,
\]
has the unique solution
\[
\boxed{
y(t)=\Phi(t,0)y_0+\int_0^t\Phi(t,s)b(s)\dd s.}
\]
Here
\[
\boxed{\Phi(t,s)=\exp\left(\int_s^t a(u)\dd u\right)}
\]
propagates the state from time $s$ to time $t$.
\end{keyproposition}

This is the integrating-factor method written so that the propagation of
initial data and forcing is explicit.

\subsubsection*{Worked forced equation}
Solve
\[
y'(t)-2y(t)=t,
\qquad y(0)=1.
\]
Here $\alpha=2$ and $g(t)=t$, hence
\[
y(t)=e^{2t}\left(1+\int_0^tse^{-2s}\dd s\right)
=\frac54e^{2t}-\frac{1+2t}{4}.
\]
Direct differentiation verifies both the ODE and the initial condition.

\subsubsection*{Worked variable-coefficient equation}
For
\[
y'(t)+\frac{2}{1+t}y(t)=1+t,
\qquad y(0)=1,
\]
the integrating factor is $M(t)=(1+t)^2$. Therefore
\[
y(t)=\frac{1}{(1+t)^2}
\left(1+\int_0^t(1+s)^3\dd s\right)
=\frac{(1+t)^4+3}{4(1+t)^2}.
\]

\subsection{Linear Systems in Several Dimensions}

The multidimensional analogue of a scalar linear equation is
\[
\boxed{x'(t)=A(t)x(t)+g(t),\qquad x(t_0)=x_0,}
\]
where $x(t)\in\R^d$, $A(t)\in\R^{d\times d}$, and $g(t)\in\R^d$. If
$A$ and $g$ are continuous, the right-hand side is continuous in $t$ and
Lipschitz in $x$ on every compact time interval. The initial value problem
therefore has a unique solution on that interval.

\subsubsection*{Homogeneous constant-coefficient systems}

For $A$ constant in time, define the matrix exponential by
\[
\boxed{e^{tA}=I+tA+\frac{t^2A^2}{2!}+\frac{t^3A^3}{3!}+\cdots.}
\]
Term-by-term differentiation yields
\[
\frac{\dd}{\dd t}e^{tA}=Ae^{tA}=e^{tA}A,
\qquad e^{0A}=I.
\]
\begin{keyproposition}{Constant-coefficient linear systems}{linear-system-solution}
For continuous $g$ and a constant matrix $A$, the initial value problem
\[
x'(t)=Ax(t)+g(t),\qquad x(0)=x_0,
\]
has the unique solution
\[
\boxed{x(t)=e^{tA}x_0+\int_0^t e^{(t-s)A}g(s)\dd s.}
\]
In particular, for the homogeneous system $g=0$,
\[
\boxed{x(t)=e^{tA}x_0.}
\]
\end{keyproposition}

This simple exponential formula relies on $A$ being constant. For a general
time-dependent $A(t)$, the expression $e^{\int_0^tA(s)\dd s}$ is not usually
valid because matrices at different times need not commute.
If $A=PDP^{-1}$ is diagonalizable, then
\[
e^{tA}=Pe^{tD}P^{-1},
\qquad
e^{tD}=\operatorname{diag}(e^{\lambda_1t},\ldots,e^{\lambda_dt}).
\]
Diagonalization therefore separates the coupled system into scalar
exponential modes. Eigenvalues with negative real parts generate decay,
positive real parts generate growth, and nonzero imaginary parts generate
oscillation.

\subsubsection*{Variation of constants in $\R^d$}

For the forced system $x'(t)=Ax(t)+g(t)$, set $x(t)=e^{tA}c(t)$. Then
$c'(t)=e^{-tA}g(t)$, which proves the formula above.
This is the multidimensional variation-of-constants, or Duhamel, formula. If
$g(t)=b$ is constant and $A$ is invertible, the equilibrium is
$x^\star=-A^{-1}b$, and
\[
\boxed{x(t)=x^\star+e^{tA}(x_0-x^\star).}
\]

\subsubsection*{Worked two-dimensional system}

Consider
\[
x'(t)=\begin{pmatrix}-1&1\\0&-2\end{pmatrix}x(t),
\qquad x(0)=\begin{pmatrix}0\\1\end{pmatrix}.
\]
The eigenpairs are
\[
\lambda_1=-1,\quad v_1=(1,0)^\top,
\qquad
\lambda_2=-2,\quad v_2=(-1,1)^\top.
\]
Since $x_0=v_1+v_2$,
\[
x(t)=e^{-t}v_1+e^{-2t}v_2
=\boxed{\begin{pmatrix}e^{-t}-e^{-2t}\\e^{-2t}\end{pmatrix}}.
\]
Both modes decay, and the slower mode $e^{-t}$ dominates at long horizons.

\subsection{Second-Order Linear Equations and First-Order Systems}

For a homogeneous constant-coefficient equation
\[
y''+ay'+by=0,
\]
try $y(t)=e^{rt}$. The characteristic equation is
\[
r^2+ar+b=0.
\]
Two distinct real roots $r_1,r_2$ give
$y=C_1e^{r_1t}+C_2e^{r_2t}$; a repeated root $r$ gives
$y=(C_1+C_2t)e^{rt}$; complex roots $\gamma\pm i\omega$ give damped
oscillations
\[
y=e^{\gamma t}(C_1\cos\omega t+C_2\sin\omega t).
\]

Every scalar second-order equation can also be written as a first-order system.
With $x_1=y$ and $x_2=y'$, the equation
\[
y''+ay'+by=f(t)
\]
becomes
\[
\boxed{
\begin{pmatrix}x_1'\\x_2'\end{pmatrix}
=\begin{pmatrix}0&1\\-b&-a\end{pmatrix}
\begin{pmatrix}x_1\\x_2\end{pmatrix}
+\begin{pmatrix}0\\f(t)\end{pmatrix}.}
\]
The characteristic roots of the scalar equation are exactly the eigenvalues
of the companion matrix.

For example,
\[
y''+3y'+2y=0,
\qquad y(0)=1,
\qquad y'(0)=0,
\]
has roots $-1$ and $-2$. The initial conditions give
\[
\boxed{y(t)=2e^{-t}-e^{-2t}}.
\]

\subsection{Numerical Methods: Explicit and Implicit Euler}

For $y'=f(t,y)$ on the grid $t_n=t_0+nh$, integrate over one time step:
\[
y(t_{n+1})=y(t_n)+\int_{t_n}^{t_{n+1}}f(s,y(s))\dd s.
\]
Approximating the integral by its left endpoint gives explicit Euler,
\[
\boxed{y_{n+1}=y_n+h f(t_n,y_n)}.
\]
Approximating it by its right endpoint gives implicit Euler,
\[
\boxed{y_{n+1}=y_n+h f(t_{n+1},y_{n+1})}.
\]

\begin{keyproposition}{First-order convergence of Euler schemes}{euler-convergence}
If $f$ is continuous in $t$, Lipschitz in $y$, and sufficiently smooth along
the exact solution, then on a fixed interval $[0,T]$ there is a constant $C_T$
independent of $h$ such that
\[
\boxed{
\max_{0\leq n\leq N}|y(t_n)-y_n|\leq C_T h.
}
\]
Thus both schemes are first-order convergent. If $f$ is $L$-Lipschitz in $y$,
the condition $hL<1$ is a simple sufficient condition for the implicit step
to define a unique $y_{n+1}$.
\end{keyproposition}

For a linear system $x'=Ax+g(t)$, the vector updates are
\[
\boxed{x_{n+1}=(I+hA)x_n+h g(t_n)}
\]
for explicit Euler, and
\[
\boxed{(I-hA)x_{n+1}=x_n+h g(t_{n+1})}
\]
for implicit Euler. Thus the implicit method requires a linear solve at each
step. If $A$ and $h$ are constant, one can factorize $I-hA$ once and reuse the
factorization throughout the time grid.

\subsubsection*{Stability on the linear test equation}

For $y'=\lambda y$, explicit Euler gives
\[
y_{n+1}=(1+h\lambda)y_n.
\]
When $\lambda<0$, numerical decay requires
\[
|1+h\lambda|<1,
\qquad\text{equivalently}\qquad
0<h<-\frac{2}{\lambda}.
\]
Implicit Euler gives
\[
y_{n+1}=\frac{1}{1-h\lambda}y_n.
\]
For every $h>0$ and every real $\lambda<0$, the update multiplier has
modulus below one. This unconditional decay stability explains its usefulness
for stiff equations.

\section{Financial Models}

\subsection{Risk-Free Growth and Discounting}

A risk-free bank account $B_t$ with a prescribed short-rate curve $r(t)$ satisfies
\[
\frac{\dd B_t}{\dd t}=r(t)B_t,
\qquad B_0>0.
\]
Solving this linear ODE gives
\[
B_t=B_0\exp\left(\int_0^t r(s)\dd s\right).
\]
For a constant rate $r$,
\[
B_t=B_0e^{rt}.
\]
Here $r(t)$ is a fixed input known as a function of time, which is why the
calculation is an ODE. If the short rate is random, its dynamics require an
SDE and pricing generally involves conditional expectations.

\subsubsection*{Continuous discounting}

The discount factor from time $0$ to time $T$ is
\[
D(0,T)=\frac{B_0}{B_T}
=\exp\left(-\int_0^T r(s)\dd s\right).
\]
For a constant rate,
\[
D(0,T)=e^{-rT}.
\]
Thus the present value of a deterministic payoff $X_T$ is
\[
X_0=D(0,T)X_T.
\]

\subsection{Mean Reversion and the Vasicek Mean Path}

The mean-reverting ODE
\[
r'(t)=a(b-r(t)),
\qquad a>0.
\]
The solution is
\[
r(t)=b+(r_0-b)e^{-at}.
\]
The integral of the rate is
\[
\int_0^T r(t)\dd t
=
bT+\frac{r_0-b}{a}(1-e^{-aT}).
\]
If this solution is used as a prescribed short-rate curve, discounting along
the mean path gives
\[
D(0,T)=
\exp\left(
-bT-\frac{r_0-b}{a}(1-e^{-aT})
\right).
\]
This is not the stochastic Vasicek zero-coupon bond price: in the stochastic
model one must take an expectation of the random discount factor.

\subsubsection*{Stochastic Vasicek model}

The full Vasicek short-rate model adds Brownian noise:
\[
\dd r_t=a(b-r_t)\dd t+\sigma\dd W_t.
\]
This is a stochastic differential equation, so a full treatment belongs to
stochastic calculus. However, ordinary calculus still explains the conditional
mean. For $t\ge s$,
\[
\frac{\dd}{\dd t}\E[r_t\mid r_s]
=a\left(b-\E[r_t\mid r_s]\right),
\]
so
\[
\E[r_t\mid r_s]
=b+(r_s-b)e^{-a(t-s)}.
\]
The parameter $a$ controls how quickly the expected short rate moves back
toward $b$, while $\sigma$ controls random fluctuations around that mean.

\end{document}
